\documentclass[a4paper,11pt]{article}

%	Tous les packages pour francisation compl\' ete
\usepackage[utf8]{inputenc}

\usepackage[T1]{fontenc}
\usepackage{lmodern,textcomp}
\usepackage[frenchb]{babel}

\usepackage{amsfonts,amsmath,amssymb,mathrsfs}
\usepackage{geometry}
\usepackage{aeguill}
\usepackage{hyperref}
\usepackage{array}
\usepackage{multirow}
\usepackage{asymptote}
\usepackage{graphicx}
\usepackage{tabularx}
\usepackage{enumerate,tablists}
\usepackage{eurosym}
\usepackage{pst-solides3d}
\usepackage{multirow}
\usepackage{extsizes}
\usepackage{calc}
\usepackage{boites}
\usepackage{color,colortbl}
\usepackage{fancyhdr}
\usepackage{stmaryrd}
\usepackage{pifont}
%\usepackage{pst-all,pstricks,pstricks-add,pst-plot,pst-eucl,pst-text,pst-tree,pst-math,pst-eps}

%\usepackage{auto-pst-pdf}
\makeatletter
\let\Test@pr@shipout\pr@shipout%% save the original definition
\let\Test@shipout\shipout
\makeatother
\usepackage{tikz,tkz-tab}
\makeatletter
\AtBeginDocument{%
  \let\pr@shipout\Test@pr@shipout%% restore it 
  \let\shipout\Test@shipout
}
\makeatother

\newtheorem{definition}{D\'efinition}
\newtheorem{theoreme}{Th\'eor\`eme}

\newcommand{\C} {\ensuremath{\mathbb{C}}}
\newcommand{\R} {\ensuremath{\mathbb{R}}}
\newcommand{\Z} {\ensuremath{\mathbb{Z}}}
\newcommand{\N} {\ensuremath{\mathbb{N}}}
\newcommand{\ds} {\displaystyle}
\newcommand{\p}{\par}
\newcommand{\abs}[1]{\left |#1\right |}
\newcommand{\eq}[1]{\underset{#1}{\sim}}
\newcommand{\tend}[1]{\underset{#1}{\longrightarrow}}
\newcommand{\norme}[1]{\left\Vert #1\right\Vert}
\newcommand{\V}[1]{\overrightarrow{#1}} %vecteur

\definecolor{gris1}{gray}{0.85}
\definecolor{gris2}{gray}{0.65}

\setlength{\textwidth} {19cm}
\setlength{\textheight} {26cm}
\addtolength{\topmargin} {-2.5cm}
\setlength{\oddsidemargin} {-1.5cm}
\setlength{\evensidemargin} {-1.5cm}

\newcommand{\encad}[1]{\begin{center}%
\fcolorbox{black}{gray!20}{\begin{minipage}[t]{0.8\linewidth}%
#1\end{minipage}}\end{center}}


\pagestyle{fancy}
\fancyfoot[R]{\today}\renewcommand\headrulewidth{1pt}
\fancyhead[L]{ESTP / Bachelor 2 / \the\year}
\fancyhead[R]{Composition de math\'ematiques}
\renewcommand\footrulewidth{1pt}
%\fancyfoot[L]{Math\' ematiques : les polyn\' emes}
%\fancyfoot[R]{\thepage}


\begin{document}
\begin{center}\bf {\Large\underline{Composition de MATH\'EMATIQUES  (3h)}}\end{center}
%\begin{center}\bf {{\small Rendre le sujet avec votre copie si le graphe de l'exercice 4 pr\' esente la courbe de la fonction $f$.}}\end{center}
\hfill\break

\underline {\ding{110}\, \bf{Exercice 1 : Intégrale } \bf{(5 points)}}
\medskip



\begin{enumerate}
    \item Montrer que l'intégrale $\displaystyle \int_{2}^{+\infty} \frac{4x}{x^4 - 1} \, dx$ converge. (2 points)
    \item Vérifier que pour tout $x \in \mathbb{R} \setminus \{-1, 1\}$ :
    \[ \frac{4x}{x^4 - 1} = -\frac{2x}{x^2 + 1} + \frac{1}{x - 1} + \frac{1}{x + 1} \] (1 point)
    \item En déduire la valeur de l'intégrale. (2 points)
\end{enumerate}



\hfill\break
\hrule
\hfill\break


\underline {\ding{110}\, \bf{Exercice 2  : Diagonalisation (12 points)}}
\medskip


\noindent On considère l'endomorphisme $f$ de $\mathbb{R}^3$ défini par :
\[ f(x, y, z) = (3x - z, 2x + 4y + 2z, -x + 3z) \]

\begin{enumerate}
    \item Déterminer la matrice $A = \mathrm{Mat}_{\mathcal{B}}(f)$ de $f$ dans la base canonique $\mathcal{B}$ de $\mathbb{R}^3$. (1 point)
    \item Déterminer le polynôme caractéristique de $f$. En déduire les valeurs propres de $f$.(3 points)
    \item Déterminer une base pour chaque espace propre de $f$. L'endomorphisme $f$ est-il diagonalisable ? (4 points)
    \item Trouver une matrice $P$ telle que $A = PDP^{-1}$, où $D$ est une matrice diagonale que l'on explicitera. (2 points)
    \item Déterminer la matrice $A^n$ pour tout $n \geq 1$.(2 points)
\end{enumerate}




\hfill\break
\hrule
\hfill\break

\underline {\ding{110}\, \bf{Exercice 3 : Dénombrement (8 points)}}\\
\medskip


Une entreprise du bâtiment TOUTBETON dispose de $10$ ouvriers spécialisés :
\begin{itemize}
    \item $4$ Maçons (numérotés de $1$ à $4$).
    \item $6$ Électriciens (numérotés de $5$ à $10$).
\end{itemize}

\begin{enumerate}
\item Pour créer un nouveau nom de dossier, la direction s'amuse à la nommer à partir d'un anagramme du nom de la société, soit ''TOUTBETON''. Combien de dossiers (d'anagrammes) existe-t-il? (2 points)
    \item 
On choisit simultanément $5$ ouvriers parmi les $10$ pour former une équipe.
\begin{enumerate}
    \item Combien y a-t-il de compositions d'équipe possibles ? (1 point)
    \item Combien d'équipes sont composées de $2$ maçons et $3$ électriciens ?(1 point)
    \item  Combien d'équipes comportent au moins un maçon ?(1 point)
\end{enumerate}

\item
On sélectionne $5$ ouvriers pour intervenir successivement sur un poste de travail (le tirage est successif et sans remise). L'ordre d'intervention est ici important.
\begin{enumerate}
    \item Combien y a-t-il de séquences d'intervention possibles ?(1 point)
    \item Combien de séquences comportent exactement $2$ maçons et $3$ électriciens ?(1 point)
    \item Combien de séquences comportent au moins un maçon ?(1 point)
\end{enumerate}


\end{enumerate}

%
%
%{\color{blue}\underline{\bf SOLUTION :}
%

%}


\end{document}
